\documentclass[sigconf,nonacm]{acmart}

\usepackage{booktabs}

% --- Strip ACM boilerplate (not part of the assignment) --- %
\setcopyright{none}
\settopmatter{printacmref=false, printccs=false, printfolios=true}
\renewcommand\footnotetextcopyrightpermission[1]{}
\acmConference[]{}{}{}
\acmISBN{}
\acmDOI{}
\acmPrice{}

% --- Placeholders: fill in before submission --- %
\newcommand{\groupnumber}{6}
\newcommand{\todo}[1]{\textcolor{red}{[#1]}}

% Configuration labels used in text, figures and tables.
\newcommand{\cfgM}{\textsf{M}}     % mutation only
\newcommand{\cfgX}{\textsf{X}}     % uniform crossover + mutation
\newcommand{\cfgH}{\textsf{H}}    % headless chicken: crossover with a random genotype + mutation
\newcommand{\cfgRS}{\textsf{RS}}   % random search

\begin{document}

% --- Cover page (not counted towards the page limit) --- %
\onecolumn
\thispagestyle{empty}
\begin{center}
  \vspace*{0.2\textheight}
  {\LARGE Evolutionary Computing\par}
  \vspace{1.5em}
  {\Large Standard Assignment -- Task II: Neuroevolution for Targeted Locomotion\par}
  \vspace{2.5em}
  {\large Group \groupnumber\par}
  \vspace{1.5em}
  \begin{tabular}{ll}
    Kimon Anagnostopoulos & 2879433 \\
    Andreas Adamantos & 2760367 \\
    Poernomo Abdul Hamid & 2746421 \\
    Youssra Ben Salah & 2769993 \\
  \end{tabular}
  \par\vfill
  {\large 13 October 2026\par}
\end{center}
\clearpage
\setcounter{page}{1}

% =========================================================================== %
%  PAPER  (max. 6 pages excluding cover page and bibliography)
% =========================================================================== %
\title{\todo{Working title:} Does Recombination Help? Crossover versus Mutation in the Neuroevolution of a Targeted-Locomotion Controller}
\author{Group \groupnumber}
\affiliation{%
  \institution{Vrije Universiteit Amsterdam}
  \city{Amsterdam}
  \country{The Netherlands}}

\begin{abstract}
\todo{Write last (4--5 sentences, as in A1): task (spider\_8, flat world, target at (3, 2) m, NN weights evolved); the four configurations and budget (9{,}030 evaluations; 15/15/10/10 runs); main numbers (\cfgX{} $0.47 \pm 0.48$ m vs.\ \cfgM{} $0.62 \pm 0.59$ m, $\hat{A}_{12} = 0.60$, n.s.; \cfgX{} beats \cfgH{} on every run; both EAs beat RS); one-line take-away (recombination is not harmful and is not mere macromutation, but its advantage over mutation is not established; remaining failures are gait speed, not steering).}
\end{abstract}

\maketitle

% --------------------------------------------------------------------------- %
\section{Introduction}
% --------------------------------------------------------------------------- %
% Rubric -- Problem definition (3 pt): research question clearly defined and stated.
% Avoid a question that is only about "how well does an EA perform".
\todo{Paragraph 1 -- context: neuroevolution of controller weights~\cite{yao1999,eiben2015}; the robot, task and framework~\cite{ariel2026}.}

\todo{Paragraph 2 -- the open question about crossover: competing conventions / permutation problem~\cite{schaffer1992,whitley1990,radcliffe1993}; mutation-only approaches~\cite{angeline1994,such2017}; evidence that the problem is overstated~\cite{hancock1992,haflidason2009}; repairs such as NEAT and neuron alignment~\cite{stanley2002,uriot2020}; early success with crossover~\cite{montana1989}. Our network has $6!\,2^6 = 46{,}080$ equivalent weight orderings per behaviour.}

\todo{Paragraph 3 -- research question (verbatim from the plan) and the headless-chicken test~\cite{jones1995} as the way to separate recombination from macromutation.}

\todo{Paragraph 4 -- hypotheses, stated before the runs. H1: \cfgX{} and \cfgM{} differ in final distance (competing conventions predicts \cfgX{} worse). H2: \cfgX{} and \cfgH{} differ (no difference would mean crossover acts as macromutation). H3: each EA reaches a lower final distance than RS.}

% --------------------------------------------------------------------------- %
\section{Methods}
% --------------------------------------------------------------------------- %
% Rubric -- Algorithm (5 pt): motivation; complete, concise description; justified
% parameters; evolution-based. Methodology (5 pt): clear procedure; everything needed
% to reproduce; appropriate plots; discussion of experimental choices.
% Rubric -- Structure: no coding details (ariel internals, SQLite, file names) here.

\begin{figure}[t]
  \centering
  \includegraphics[width=\columnwidth]{figures/fig1_convergence.pdf}
  \caption{Search dynamics over 300 generations (9{,}030 evaluations, top axis). Solid:
  best distance to the target found so far; dashed: mean distance of the population
  (EA variants only). Curves show the mean over independent runs (15 for \cfgM{} and \cfgX{},
  10 for \cfgH{} and \cfgRS{}) and the shaded bands $\pm 1$ standard deviation over runs.
  Lower is better.}
  \label{fig:convergence}
\end{figure}

\subsection{Task, robot and fitness}
\todo{spider\_8 body from the John Set (8 hinges) on flat ground; spawn at the origin, target at $(3, 2)$~m ($\approx 3.6$~m, off-axis so the controller has to steer); 15~s of simulated time at a 2~ms physics step. Fitness $f = \lVert \mathbf{p}_{xy}(T) - \mathbf{t}_{xy} \rVert$, minimised. Justify body and target with the pilot (Table~\ref{tab:pilots}): the 2~m straight-ahead target was nearly solved within 520 evaluations, which would make the arms indistinguishable.}

\subsection{Controller and representation}
\todo{Feed-forward network, 14 inputs $\to$ 6 tanh $\to$ 8 tanh, outputs scaled to $\pm\pi/2$ and applied directly as hinge targets at every physics step. Inputs: 8 hinge angles (scaled), unit vector to the target in the robot's heading frame (2), clipped distance (1), $\sin$ and $\cos$ of a fixed 1~Hz clock (2), bias (1). Argue why: target direction for steering, clock for rhythm; the clock has no evolved parameters, so this is not a CPG. Genotype: the 132 weights as a real vector; initial weights $\mathcal{N}(0, 0.5^2)$.}

\subsection{Evolutionary algorithm and variation}
\label{sec:ea}
\todo{$(\mu+\lambda)$, $\mu = \lambda = 30$, tournament $k = 3$, truncation survivor selection (elitist), 300 generations. The four configurations differ only in how a child is formed before the identical Gaussian mutation ($\sigma = 0.2$, per-gene $p = 0.1$, redrawn until at least one gene changes, so that no evaluation is spent on a clone): \cfgM{} copies one parent; \cfgX{} applies uniform crossover to two tournament winners; \cfgH{} applies the same crossover to one tournament winner and a freshly drawn random genotype that is never evaluated~\cite{jones1995}; \cfgRS{} samples genotypes from the initialiser. Motivate uniform crossover (the standard operator for real vectors, maximally disruptive to hidden units).}

\begin{table}[t]
  \caption{Parameter settings shared by all EA variants.}
  \label{tab:params}
  \small
  \centering
  \begin{tabular}{@{}ll@{}}
    \toprule
    Parameter & Value \\
    \midrule
    Population size $\mu$ / offspring $\lambda$ & 30 / 30 \\
    Generations & 300 \\
    Evaluation budget per run & $\mu + 300\lambda = 9{,}030$ \\
    Parent selection & tournament, $k = 3$, with replacement \\
    Survivor selection & $(\mu+\lambda)$ truncation \\
    Crossover (\cfgX{}, \cfgH{}) & uniform, swap probability 0.5 \\
    Mutation & Gaussian, $\sigma = 0.2$, per-gene $p = 0.1$ \\
    Initial weights & $\mathcal{N}(0, 0.5^2)$, 132 genes \\
    Network & 14 -- 6 -- 8, tanh \\
    Simulated time per evaluation & 15~s (2~ms steps) \\
    Independent runs & 15 (\cfgM{}, \cfgX{}), 10 (\cfgH{}, \cfgRS{}) \\
    \bottomrule
  \end{tabular}
\end{table}

\subsection{Pilot experiments}
\todo{Pilot seeds (100--102) never reused. Body pilot: short \cfgM{} run on all 13 John Set bodies. $\sigma$ pilot: \cfgM{} and \cfgX{} at 1{,}530 evaluations, 3 seeds, $\sigma \in \{0.05, 0.1, 0.2, 0.4\}$. Report the deviation honestly: the planned rule (best mean over both arms) picks 0.05 by 0.015~m, a tie at $n = 3$ that is best for \cfgX{} and worst for \cfgM{}; we chose 0.2 instead, the best worst-arm value, so that neither arm is favoured. Budget raised from 150 to 300 generations because no arm had reached the plateau criterion at 150.}

\begin{table}[t]
  \caption{Pilot experiments. Left: the four best John Set bodies in a short mutation-only
  run (520 evaluations, 2 seeds, 2~m straight-ahead target; mean final distance in m).
  Right: mean final distance (m) of \cfgM{} and \cfgX{} after 1{,}530 evaluations (3 seeds)
  for each mutation step size, and of the worse of the two arms.}
  \label{tab:pilots}
  \small
  \centering
  \begin{tabular}{@{}lc@{\hspace{1.2em}}lccc@{}}
\toprule
Body & Best & $\sigma$ & M & X & Worse \\
\midrule
\texttt{linkin\_modified} & 0.43 & 0.05 & 1.88 & 1.11 & 1.88 \\
\texttt{spider\_8} & 0.59 & 0.1 & 1.50 & 1.55 & 1.55 \\
\texttt{centipede\_4} & 0.90 & 0.2 & 1.53 & 1.50 & 1.53 \\
\texttt{turtle} & 1.09 & 0.4 & 1.59 & 1.67 & 1.67 \\
\bottomrule
\end{tabular}

\end{table}

\subsection{Experimental setup and statistics}
\todo{Seeds 0--14 for \cfgM{} and \cfgX{}, 0--9 for \cfgH{} and \cfgRS{}; for a given seed every configuration starts from the same initial population, so runs are paired. Primary test: two-sided Mann--Whitney $U$~\cite{mann1947} on final best distance, Holm-corrected~\cite{holm1979} over the five comparisons of Table~\ref{tab:results}, $\alpha = 0.05$; effect size $\hat{A}_{12}$~\cite{vargha2000} with bootstrap 95\% CI and a pre-registered smallest effect of interest of $|\hat{A}_{12} - 0.5| \geq 0.21$; paired Wilcoxon~\cite{wilcoxon1945} as a secondary check. Exploratory, outside the Holm family: distance at 20\% of the budget, and the fraction of runs ending within 0.1~m of the target (Fisher's exact test). Plateau criterion: mean best-so-far improves by less than 1\% over the last 20\% of evaluations. Reproducibility: runs are deterministic given the seed; code, commands, per-run configuration (seed, parameters, commit, package versions) are included.}

% --------------------------------------------------------------------------- %
\section{Results and discussion}
% --------------------------------------------------------------------------- %
% Rubric -- Results & Discussion (10 pt): results presented clearly with tables/plots;
% analysed thoroughly; linked back to the research question; strong critical
% reflection (not just numbers); limitations and unexpected findings discussed.

\begin{table}[t]
  \caption{Final best distance to the target in m (lower is better) and the comparisons
  of the hypotheses. Reached: runs ending within 0.1~m of the target. $\hat{A}_{12}$ is the
  probability that the first configuration of a pair attains the lower value; Wins counts
  seeds on which it does (paired by seed). $p_\text{Holm}$: two-sided Mann--Whitney $U$,
  Holm-corrected over the five comparisons.}
  \label{tab:results}
  \small
  \centering
  \begin{tabular}{@{}lcccc@{}}
\toprule
Config. & Mean $\pm$ SD & Median [IQR] & Range & Reached \\
\midrule
M & $0.62 \pm 0.59$ & 0.46 [0.98] & 0.002--1.91 & 5/15 \\
X & $0.47 \pm 0.48$ & 0.39 [0.90] & 0.001--1.18 & 7/15 \\
H & $2.27 \pm 0.21$ & 2.32 [0.31] & 1.881--2.54 & 0/10 \\
RS & $2.42 \pm 0.13$ & 2.46 [0.20] & 2.209--2.59 & 0/10 \\
\midrule
Comparison & $p_\text{Holm}$ & $\hat{A}_{12}$ [95\% CI] & Wins & Hyp. \\
\midrule
X vs.\ M & 0.384 & 0.60 [0.38, 0.79] & 8/15 & H1 \\
X vs.\ H & $1.8\times10^{-4}$ & 1.00 [1.00, 1.00] & 10/10 & H2 \\
M vs.\ RS & $1.8\times10^{-4}$ & 1.00 [1.00, 1.00] & 10/10 & H3 \\
X vs.\ RS & $1.8\times10^{-4}$ & 1.00 [1.00, 1.00] & 10/10 & H3 \\
H vs.\ RS & 0.178 & 0.73 [0.49, 0.92] & 6/10 & H3 \\
\bottomrule
\end{tabular}

\end{table}

\begin{figure}[t]
  \centering
  \includegraphics[width=\columnwidth]{figures/fig2_final_best.pdf}
  \caption{Final best distance to the target of every independent run (dots) after
  9{,}030 evaluations. Lower is better.}
  \label{fig:final}
\end{figure}

\begin{figure*}[t]
  \centering
  \includegraphics[width=\textwidth]{figures/fig3_offspring.pdf}
  \caption{Left: fraction of runs whose best controller is within 0.1~m of the target.
  Centre: mean pairwise Euclidean distance between the genotypes of the population.
  Right: fraction of children better than the better of their evaluated parents
  (5-generation moving average, log scale). Means over runs.}
  \label{fig:offspring}
\end{figure*}

\begin{figure*}[t]
  \centering
  \includegraphics[width=\textwidth]{figures/fig4_paths.pdf}
  \caption{Top-down trajectories of the core of the best controller of every run, from
  the start ($\blacksquare$) towards the target ($\times$, circle of radius 0.1~m). Dots
  mark the final positions.}
  \label{fig:paths}
\end{figure*}

\begin{figure}[t]
  \centering
  \includegraphics[width=\columnwidth]{figures/fig5_paired.pdf}
  \caption{Final best distance of \cfgX{} against \cfgM{} for the same seed, i.e.\ from the same
  initial population (log scales). Points below the diagonal are seeds on which crossover
  did better; dotted lines mark 0.1~m.}
  \label{fig:paired}
\end{figure}

\paragraph{Crossover versus mutation (H1).}
\todo{\cfgX{} lower on mean, median and worst case ($0.47 \pm 0.48$ vs.\ $0.62 \pm 0.59$ m; worst 1.18 vs.\ 1.91), more runs reach the target (7/15 vs.\ 5/15) and earlier (median $\approx$3{,}870 vs.\ 5{,}820 evaluations; Fig.~\ref{fig:offspring} left), but $p_\text{Holm} = 0.38$, paired Wilcoxon $p = 0.52$, $\hat{A}_{12} = 0.60$ [0.38, 0.79]: the CI covers both no effect and the smallest effect of interest, so H1 is inconclusive -- neither supported nor refuted. Per seed (Fig.~\ref{fig:paired}) \cfgX{} wins 8/15, with two seeds where only \cfgX{} reaches the target and two where only \cfgM{} does. Against the competing-conventions prediction, crossover did not hurt.}

\paragraph{Dynamics and diversity.}
\todo{\cfgM{} faster for the first 40 generations, \cfgX{} overtakes at generation 40 ($\approx$1{,}230 evaluations) and stays ahead to the end (Fig.~\ref{fig:convergence}); \cfgM{} diversity falls below 2 by generation 11, \cfgX{} only by generation 54, and settles higher (1.7 vs.\ 1.2; Fig.~\ref{fig:offspring} centre). Interpretation: crossover first slows the climb and then explores more; once the population has converged, crossing similar parents is a small step, which is why the arms look alike late. Exploratory early comparison at 20\% of the budget: $\hat{A}_{12} = 0.65$, n.s.}

\paragraph{Recombination versus macromutation (H2).}
\todo{\cfgX{} beats \cfgH{} on every run ($\hat{A}_{12} = 1.00$, $p_\text{Holm} = 1.8\times10^{-4}$): H2 supported -- combining two selected parents is not a random macromutation. Critical point: with uniform crossover the random donor replaces half the weights, so \cfgH{} is close to random search (diversity stays $\approx 8$; children rarely improve, Fig.~\ref{fig:offspring} right); the control therefore shows that recombination is not a large random jump, not that it transfers building blocks.}

\paragraph{Comparison with random search (H3).}
\todo{\cfgM{} and \cfgX{} beat \cfgRS{} on every seed ($p_\text{Holm} = 1.8\times10^{-4}$). \cfgH{} vs.\ \cfgRS{}: $\hat{A}_{12} = 0.73$ but $p_\text{Holm} = 0.18$ -- H3 holds for the EAs with a selected-parent variation only.}

\paragraph{What the controllers do.}
\todo{Fig.~\ref{fig:paths}: every EA run walks a nearly straight line towards the off-axis target, so steering is solved in all runs; unsuccessful runs stop short on the right line -- the remaining error is gait speed, not navigation, which explains the split into runs that reach the target and runs that stall at 0.3--1.9~m. \cfgRS{} and \cfgH{} controllers head roughly the right way but end only 1.1--1.8~m from the start (EAs: 1.7--3.7~m). The core never dropped below its resting height (0.074--0.077~m) in any best run; this rules out collapse, not flipping, so check the videos before claiming more.}

\paragraph{Limitations.}
\todo{(1) Plateau criterion not met for the EAs (mean improvement 7.3\% and 6.1\% over the last 20\%, against 1\%; halved from 150 generations) -- longer runs could change the gap. (2) Power: 15 runs per arm and the bimodal outcome leave H1 open. (3) One shared $\sigma$, chosen by a minimax rule; \cfgX{} preferred a smaller $\sigma$ in the pilot. (4) Uniform crossover mixes generic disruption with permutation effects, so the design does not isolate competing conventions. (5) One body, one world, one target, 15~s episodes.}

% --------------------------------------------------------------------------- %
\section{Conclusions}
% --------------------------------------------------------------------------- %
% Rubric -- Conclusion (5 pt): strong and concise; summarises main findings; directly
% answers the research question; limitations and future work.
\todo{Answer the RQ directly: recombination did not hurt and is clearly not a macromutation (H2), but its advantage over mutation alone was not established at this budget and sample size (H1 inconclusive); both EAs beat random search (H3). Main insight: crossover delays convergence and keeps diversity; failures are about speed, not steering. Future work: unit-wise crossover that moves whole hidden neurons, neuron-aligned crossover~\cite{uriot2020}, each arm at its own step size, more runs and longer budgets.}

\bibliographystyle{ACM-Reference-Format}
\bibliography{references}

\end{document}
